\documentclass[10pt]{amsart}
\usepackage[a4paper,margin=22mm]{geometry}
\usepackage[no-math]{fontspec}
\setmainfont{Latin Modern Roman}[SmallCapsFont={lmromancaps10-regular.otf}]
\usepackage{amsmath,amssymb,amsthm,mathtools,mathrsfs,hyperref,graphicx,etoolbox}
\emergencystretch=3em
\allowdisplaybreaks
\newtheorem{theo}{Theorem}[section]
\newtheorem{lemm}[theo]{Lemma}
\newtheorem{prop}[theo]{Proposition}
\theoremstyle{definition}
\newtheorem{defi}[theo]{Definition}
\newtheorem{rema}[theo]{Remark}
\newtheorem{exam}[theo]{Example}
\numberwithin{equation}{section}
\newcommand{\bbR}{\mathbb{R}}
\newcommand{\bbN}{\mathbb{N}}
\newcommand{\mfs}{\mathfrak{s}}
\newcommand{\wco}{w^*}
\newcommand{\vco}{v^*}

\newcommand{\bfa}{\mathbf{a}}
\newcommand{\bfb}{\mathbf{b}}
\newcommand{\bfc}{\mathbf{c}}
\newcommand{\bfzo}{\mathbf{0}}
\newcommand{\bfaob}{\mathbf{a}\ominus\mathbf{b}}
\newcommand{\bfcoa}{\mathbf{c}\ominus\mathbf{a}}
\newcommand{\bfcpb}{\mathbf{c}\oplus\beta}

\newcommand{\bfA}{\mathbf{A}}
\newcommand{\bfT}{\mathbf{T}}
\newcommand{\bfG}{\mathbf{G}}
\newcommand{\bfk}{\mathbf{k}}
\newcommand{\bfl}{\mathbf{l}}

\newcommand{\mcL}{\mathcal{L}}
\newcommand{\mcD}{\mathcal{D}}
\newcommand{\mcR}{\mathcal{R}}
\newcommand{\mcK}{\mathcal{K}}
\newcommand{\mcI}{\mathcal{I}}
\newcommand{\mcJ}{\mathcal{J}}
\newcommand{\mcN}{\mathcal{N}}
\newcommand{\mcA}{\mathcal{A}}
\newcommand{\mcB}{\mathcal{B}}
\newcommand{\mcC}{\mathcal{C}}
\newcommand{\scT}{\mathscr{T}}
\newcommand{\scM}{\mathscr{M}}

%\newcommand{\id}{\mathop{\text{\rm id}}}
\DeclareMathOperator{\id}{id}
\DeclareMathOperator{\Rre}{Re}
\DeclareMathOperator{\spa}{span}
\DeclareMathOperator{\deff}{def}
\newcommand{\tri}{|\!|\!|}
\newcommand{\llparenthesis}{\mathopen{(\mkern -2.7mu|}}
\newcommand{\rrparenthesis}{\mathclose{|\mkern -2.7mu)}}
\newcommand{\llbracket}{\mathopen{[\mkern -2.7mu[}}
\newcommand{\rrbracket}{\mathclose{]\mkern -2.7mu]}}
\newcommand{\llbigparenthesis}{\mathopen{\big(\mkern -2.7mu\big|}}
\newcommand{\rrbigparenthesis}{\mathclose{\big|\mkern -2.7mu\big)}}
\newcommand{\lbp}{\llbigparenthesis}
\newcommand{\rbp}{\rrbigparenthesis}
\newcommand{\lp}{\llparenthesis}
\newcommand{\rp}{\rrparenthesis}
\newcommand{\lb}{\llbracket}
\newcommand{\rb}{\rrbracket}

\newcommand{\bigtri}{\big|\mkern -2.7mu\big|\mkern -2.7mu\big|}
\newcommand{\llbigbracket}{\mathopen{\big[\mkern -4.2mu\big[}}
\newcommand{\rrbigbracket}{\mathclose{\big]\mkern -4.2mu\big]}}
\newcommand{\lbb}{\llbigbracket}
\newcommand{\rbb}{\rrbigbracket}

\newenvironment{alcases}{\left\{\begingroup\arraycolsep=0pt\renewcommand{\arraystretch}{1.2}\begin{array}{r@{\;}l@{\quad}l}}{\end{array}\endgroup\right.}


\newenvironment{step}[1]{\refstepcounter{stepcount}\begin{proof}[{\bf\thestepcount[#1]}}{\let\qed\relax\end{proof}}

%%%%%%%%%%%%%%%%%%%%%%%%%%%%%%%%%%%%%%%%%%%%%%%%%%%%%%%%%%
\graphicspath{{./figures/}}

\newcommand*{\mk}{\mkern -1mu}
\newcommand*{\Mk}{\mkern -2mu}
\newcommand*{\mK}{\mkern 1mu}
\newcommand*{\MK}{\mkern 2mu}

\hypersetup{urlcolor=purple, linkcolor=blue, citecolor=red}

\newcommand*{\relabel}{\renewcommand{\labelenumi}{(\theenumi)}}
\newcommand*{\romanenumi}{\renewcommand*{\theenumi}{\roman{enumi}}\relabel}
\newcommand*{\Romanenumi}{\renewcommand*{\theenumi}{\Roman{enumi}}\relabel}
\newcommand*{\alphenumi}{\renewcommand*{\theenumi}{\alph{enumi}}\relabel}
\newcommand*{\Alphenumi}{\renewcommand*{\theenumi}{\Alph{enumi}}\relabel}
\title[Single-model weighted Besov reconstruction]{Single-model weighted Besov reconstruction for regularity-integrability structures under a general G-type semigroup}
\author{Masato Hoshino}
\date{Annales Henri Lebesgue 8 (2025),151–180; DOI10.5802/ahl.232}
\hypersetup{hidelinks}
\AtBeginDocument{\let\originalRef\ref\renewcommand{\ref}[1]{\ifstrequal{#1}{section:schauder}{\href{https://doi.org/10.5802/ahl.232}{5 (source)}}{\originalRef{#1}}}}
\begin{document}
\maketitle
\setcounter{section}{1}
\section{Besov spaces associated with the semigroup of operators}\label{section:semigroup}

In this section, we define the Besov norms associated with the semigroup of operators. For the sake of generality, we define the weighted Besov norms with arbitrary integrability exponents $p,q\in[1,\infty]$.

\subsection{Notations}\label{subsection:notation}

The symbol $\bbN$ denotes the set of all nonnegative integers. Throughout this paper, we fix an integer $d\ge1$, the \emph{scaling} $\mfs=(\mfs_1,\dots,\mfs_d)\in[1,\infty)^d$, and a number $\ell>0$. We define $|\mfs|=\sum_{i=1}^d\mfs_i$. For any multiindex ${\bfk}=(k_i)_{i=1}^d\in\bbN^d$, any $x=(x_i)_{i=1}^d\in\bbR^d$, and any $t>0$, we define the notations
\begin{align*}
&{\bfk}!:=\prod_{i=1}^dk_i!,\quad |{\bfk}|_{\mfs}:=\sum_{i=1}^d\mfs_ik_i,\quad
\|x\|_\mfs:=\sum_{i=1}^d|x_i|^{1/\mfs_i},\\
&x^{\bfk}:=\prod_{i=1}^dx_i^{k_i},\quad t^{\mfs/\ell}x:=\left(t^{\mfs_i/\ell}x_i\right)_{i=1}^d,\quad t^{-\mfs/\ell}x:=\left(t^{-\mfs_i/\ell}x_i\right)_{i=1}^d.
\end{align*}
We define the set $\bbN[\mfs]:=\{|{\bfk}|_\mfs\,;\,{\bfk}\in\bbN^d\}$, which will be used in Section~\ref{section:schauder}. The parameter $t$ is not a physical time variable, but an auxiliary variable used to define regularities of distributions. For multiindices ${\bfk}=(k_i)_{i=1}^d$ and ${\bfl}=(l_i)_{i=1}^d$, we write ${\bfl}\le{\bfk}$ if $l_i\le k_i$ for any $1\le i\le d$, and then define $\binom{{\bfk}}{{\bfl}}:=\prod_{i=1}^d\binom{k_i}{l_i}$.

We also fix a nonnegative measurable function $G:\bbR^d\to\bbR$ and define for any $t>0$,
\[
G_t(x)=t^{-|\mfs|/\ell}G\big(t^{-\mfs/\ell}x\big).
\]

We use the notation $A\lesssim B$ for two functions $A(x)$ and $B(x)$ of a variable $x$, if there exists a constant $c>0$ independent of $x$ such that $A(x)\le cB(x)$ for any $x$.

\subsection{Weighted \texorpdfstring{$L^p$}{Lp} spaces}

First we introduce the class of appropriate weight functions.

\begin{defi}\label{defweight}
A continuous function $w:\bbR^d\to(0,1]$ is called a \emph{weight}. A weight $w$ is said to be \emph{$G$-controlled} if there exists a continuous function $\wco:\bbR^d\to[1,\infty)$ such that
\begin{align}\label{weightmoderate}
w(x+y)\le \wco(x)w(y)
\end{align}
for any $x,y\in\bbR^d$ and
\begin{align}\label{weightintegrable}
\sup_{0\,<\,t\,\le\,T}\;\sup_{x\,\in\,\bbR^d}\Big\{\|x\|_\mfs^n\,\wco\left(t^{\mfs/\ell}x\right)G(x)\Big\}<\infty
\end{align}
for any $n\ge0$ and $T>0$. (In the terminology of~\cite[Definition~2.3]{MW17}, $w$ is said to be \emph{$\wco$-moderate}.)
\end{defi}

\begin{defi}
For any $p\in[1,\infty]$ and any weight $w$, we define the weighted $L^p$ norm of a measurable function $f:\bbR^d\to\bbR$ by
\begin{align*}
\|f\|_{L^p(w)}:=\|fw\|_{L^p(\bbR^d)}.
\end{align*}
We denote by $L^p(w)$ the space of all measurable functions with finite $L^p(w)$ norms, and define the closed subspace $L_c^p(w)$ as the completion of the set $C(\bbR^d)\cap L^p(w)$ under the $L^p(w)$ norm.
\end{defi}

Note that $\|\cdot\|_{L^p(w)}$ is nondegenerate because $w$ is fully supported in $\bbR^d$ by definition. An advantage of introducing $L_c^p(w)$ is that we can use density arguments. Although the space $L_c^\infty(w)$ is strictly smaller than $L^\infty(w)$, it is often sufficient to consider $L_c^\infty(w)$ in applications to SPDEs. In the following, we prove three useful inequalities.

\begin{lemm}\label{lem:weightedholder}
Let $p,q,r\in[1,\infty]$ be such that $1/r=1/p+1/q$. For any weights $w_1,w_2$ and functions $f\in L^p(w_1)$ and $g\in L^q(w_2)$, we have
\[
\|fg\|_{L^r(w_1w_2)}\le\|f\|_{L^p(w_1)}\|g\|_{L^q(w_2)}.
\]
\end{lemm}

\begin{proof}
Since $\|fg\|_{L^r(w_1w_2)}=\|(fw_1)(gw_2)\|_{L^r}$, the result follows from H\"older's inequality.
\end{proof}

\begin{lemm}[{\cite[Theorem~2.4]{MW17}}]\label{lem:weightedyoung}
Let $w$ be a $G$-controlled weight. For any $T>0$, there exists a constant $C_T>0$ depending only on $G,\wco$, and $T$, such that, for any $t\in(0,T]$, $1\le p\le q\le\infty$, and $f\in L^p(w)$, we have
\begin{align*}
\|G_t*f\|_{L^q(w)}\le C_T\,t^{-\frac{|\mfs|}\ell\left(\frac1p-\frac1q\right)}\|f\|_{L^p(w)}.
\end{align*}
\end{lemm}

\begin{proof}
Since $|(G_t*f)(x)w(x)|\le((G_t\wco)*(|f|w))(x)$ by the inequality~\eqref{weightmoderate}, the result follows from Young's inequality. The proportional constant is $\|G_t\wco\|_{L^r(\bbR^d)}$, where $\frac1r=1+\frac1q-\frac1p$. Since
\begin{align*}
\left\|G_t\wco\right\|_{L^r(\bbR^d)}&=\left\|t^{-|\mfs|/\ell}G\left(t^{-\mfs/\ell}x\right)\wco(x)\right\|_{L_x^r(\bbR^d)}\\
&=t^{-\frac{|\mfs|}\ell\left(1-\frac1r\right)}\left\|G(x)\wco\left(t^{\mfs/\ell}x\right)\right\|_{L_x^r(\bbR^d)}\\
& =t^{-\frac{|\mfs|}\ell\left(\frac1p-\frac1q\right)}\left\|G(x)\wco\left(t^{\mfs/\ell}x\right)\right\|_{L_x^r(\bbR^d)}
\end{align*}
by the scaling property, we have the result by using the condition~\eqref{weightintegrable}.
\end{proof}



\begin{lemm}\label{lem:weightedtranslation}
Let $w$ be a $G$-controlled weight. For any $p\in[1,\infty]$, $f\in L^p(w)$, and $h\in\bbR^d$, we have
\begin{align*}
\|f(\cdot-h)\|_{L^p(w)}\le \wco(h)\|f\|_{L^p(w)}.
\end{align*}
\end{lemm}

\begin{proof}
The result follows from the inequality $|f(x-h)|w(x)\le \wco(h)|f(x-h)|w(x-h)$ and the translation invariance of the unweighted $L^p(\bbR^d)$ norm.
\end{proof}

\subsection{Semigroup of operators}

We introduce a semigroup of operators.

\begin{defi}\label{asmp1}
We call a family of continuous functions $\{Q_t:\bbR^d\times\bbR^d\to\bbR\}_{t\,>\,0}$ a \emph{$G$-type semigroup (of operators)} if it satisfies the following properties.
\begin{enumerate}\romanenumi
%\renewcommand{\theenumi}{(\roman{enumi})}
%\renewcommand{\labelenumi}{(\roman{enumi})}
\item\label{asmp1:semigroup}
(Semigroup property) For any $0<s<t$ and $x,y\in\bbR^d$,
\[
\int_{\bbR^d}Q_{t-s}(x,z)Q_s(z,y)dz=Q_t(x,y).
\]
\item\label{asmp1:conservative}
(Conservativity) For any $x\in\bbR^d$,
\[
\lim_{t\downarrow0}\int_{\bbR^d}Q_t(x,y)dy=1.
\]
\item\label{asmp1:gauss}
(Upper $G$-type estimate) There exists a constant $C_1>0$ such that, for any $t>0$ and $x,y\in\bbR^d$,
\[
|Q_t(x,y)|\le C_1G_t(x-y).
\]
\item\label{asmp2:derivative}
(Time derivative) For any $x,y\in\bbR^d$, $Q_t(x,y)$ is differentiable with respect to $t$. Moreover, there exists a constant $C_2>0$ such that, for any $t>0$ and $x,y\in\bbR^d$,
\[
|\partial_tQ_t(x,y)|\le C_2\, t^{-1}G_t(x-y).
\]
\end{enumerate}
\end{defi}

\stepcounter{theo}
We identify the function $Q_t(x,y)$ with a continuous linear operator on $L^p(w)$ by
\[
(Q_tf)(x):=:Q_t(x,f):=\int_{\bbR^d}Q_t(x,y)f(y)dy,\qquad f\in L^p(w),\ x\in\bbR^d.
\]
Note that $Q_t$ is closed in $L^p(w)$ because of Definition~\ref{asmp1}$\MK$\eqref{asmp1:gauss} and Lemma~\ref{lem:weightedyoung}.

\begin{prop}\label{prop:operatorQt}
Let $w$ be a $G$-controlled weight and let $p\in[1,\infty]$. For any $f\in L^p(w)$ and $t>0$, $Q_tf$ is a continuous function. In addition, if $f\in C(\bbR^d)\cap L^p(w)$, then
\[
\lim_{t\downarrow0}(Q_tf)(x)=f(x)
\]
for any $x\in\bbR^d$.
\end{prop}

\begin{proof}
Let $f\in L^p(w)$. To show the continuity of $(Q_tf)(x)$ with respect to $x$, it is sufficient to consider the case $t=1$ and $x=0$. Note that, in the region $\|x\|_\mfs\le1$, we have
\begin{align*}
|Q_1(x,y)f(y)w(x)| &\lesssim |G(x-y)\wco(x-y)||f(y)w(y)|\lesssim \frac{|f(y)w(y)|}{1+\|y\|_\mfs^n}
\end{align*}
for any $n\ge0$ by the property~\eqref{weightintegrable}. This implies that
\begin{align*}
\lim_{x\to0}(Q_1f)(x)w(x)=\int_{\bbR^d}\lim_{x\to0}Q_1(x,y)f(y)w(x)dy=(Q_1f)(0)w(0)
\end{align*}
by Lebesgue's convergence theorem. Since $w$ is strictly positive and continuous, we have $\lim_{x\to0}(Q_1f)(x)=(Q_1f)(0)$.

Next let $f\in C(\bbR^d)\cap L^p(w)$. To show the continuity with respect to $t$, it is sufficient to consider the case $x=0$. For any $\varepsilon>0$, we can choose $\delta>0$ such that $|f(y)-f(0)|<\varepsilon$ for any $\|y\|_\mfs<\delta$, and have
\begin{multline*}
\big|w(0)(Q_tf-f)(0)\big|\\
\begin{aligned}
&=w(0)\left|\int_{\bbR^d}Q_t(0,y)\big(f(y)-f(0)\big)dy+\left(\int_{\bbR^d}Q_t(0,y)dy-1\right)f(0)\right|\\
&\le w(0)\varepsilon\int_{\|y\|_\mfs\,<\,\delta}G_t(-y)dy+w(0)\int_{\|y\|_\mfs\,\ge\,\delta}G_t(-y)|f(y)|dy\\
&\quad+w(0)\int_{\|y\|_\mfs\,\ge\,\delta}G_t(-y)|f(0)|dy+w(0)|f(0)|\left|\int_{\bbR^d}Q_t(0,y)dy-1\right|.
\end{aligned}
\end{multline*}
In the far right-hand side, the only nontrivial part is the second term. We bound it from above by
\[
\int_{\|y\|_\mfs\,\ge\,\delta}G_t(-y)\wco(-y)|f(y)|w(y)dy\le\|(G_t\wco)(y)\|_{L^{p'}(\|y\|_\mfs\,\ge\,\delta)}\|fw\|_{L^p\left(\bbR^d\right)},
\]
where $1/p+1/{p'}=1$. We then have that $\|(G_t\wco)(y)\|_{L^{p'}(\|y\|_\mfs\,\ge\,\delta)}\to0$ as $t\downarrow0$ by the property~\eqref{weightintegrable}.
\end{proof}

\subsection{Besov spaces associated with semigroup}

In what follows, we fix a $G$-type semigroup $\{Q_t\}_{t\,>\,0}$ and a $G$-controlled weight $w$. We define the weighted Besov spaces associated with $\{Q_t\}_{t\,>\,0}$, as studied in~\cite{BB16,BDY12, GL15}.

\begin{defi}
For every $\alpha\le0$ and $p,q\in[1,\infty]$, we define the Besov space $B_{p,q}^{\alpha,Q}(w)$ as the completion of $L_c^p(w)$ under the norm
\[
\|f\|_{B_{p,q}^{\alpha,Q}(w)}:=\|Q_1f\|_{L^p(w)}+\big\|t^{-\alpha/\ell}\|Q_t f\|_{L^p(w)}\big\|_{L^q\left(0,1;\,t^{-1}dt\right)}.
\]
\end{defi}

When $\mfs=(1,1,\dots,1)$, $\ell=2$, and $Q_t$ is the heat semigroup $e^{t\Delta}$, the above norm (with $\alpha<0$ and $w=1$) is equivalent to the classical Besov norm in Euclidean setting. See e.g., \cite[Theorem~2.34]{BCD11} or~\cite[Theorem~2.6.4]{Tri92}.

\begin{rema}\label{QtsymmetricS'}
We can see that $\|\cdot\|_{B_{p,q}^{\alpha,Q}(w)}$ is nondegenerate in $L_c^p(w)$ by the temporal continuity of $Q_t$ (Proposition~\ref{prop:operatorQt}) and the density argument. This is the only reason why we define the Besov spaces from $L_c^p(w)$ as above. On the other hand, if $Q_t$ is symmetric in the sense that $Q_t(y,x)=Q_t(x,y)$ for any $x,y\in\bbR^d$, then for any locally integrable function $f$ and $\varphi\in C_0^\infty(\bbR^d)$, we have
\[
\int_{\bbR^d}(Q_tf)(x)\varphi(x)dx=\int_{\bbR^d}f(x)(Q_t\varphi)(x)dx\xrightarrow{t\downarrow0} \int_{\bbR^d}f(x)\varphi(x)dx,
\]
which implies $Q_tf\to f$ as $t\downarrow0$ in the distributional sense. In this case, $\|\cdot\|_{B_{p,q}^{\alpha,Q}(w)}$ is nondegenerate in whole $L^p(w)$.
\end{rema}

The following result implies that we can ignore the difference of the parameter $q$ at the cost of infinitesimal difference of the parameter $\alpha$. Therefore, we pay less attention to $q$ in this paper and write $B_p^{\alpha,Q}(w):=B_{p,\infty}^{\alpha,Q}(w)$.

\begin{prop}\label{prop:almostequiv}
For any $f\in L^p(w)$, the following inequalities hold.
\begin{enumerate}
%\renewcommand{\theenumi}{(\arabic{enumi})}
%\renewcommand{\labelenumi}{(\arabic{enumi})}
\item\label{almostequiv2}
For any $1\le q_1\le q_2\le\infty$ and $\alpha_1<\alpha_2\le0$, we have $\|f\|_{B_{p,q_1}^{\alpha_1,Q}(w)}\lesssim\|f\|_{B_{p,q_2}^{\alpha_2,Q}(w)}$.
\item\label{almostequiv3}
For any $\alpha\le0$, we have $\|f\|_{B_{p,\infty}^{\alpha,Q}(w)}\lesssim\|f\|_{B_{p,1}^{\alpha,Q}(w)}$.
\end{enumerate}
Here the implicit proportional constants depend only on $G,\wco$, and the regularity and integrability exponents.
\end{prop}

\begin{rema}
As a result of~\eqref{almostequiv2} and~\eqref{almostequiv3}, we have
\begin{align*}
\|f\|_{B_{p,q_1}^{\alpha-2\varepsilon,Q}(w)}
\lesssim\|f\|_{B_{p,\infty}^{\alpha-\varepsilon,Q}(w)}
\lesssim\|f\|_{B_{p,1}^{\alpha-\varepsilon,Q}(w)}
\lesssim\|f\|_{B_{p,q_2}^{\alpha,Q}(w)}
\end{align*}
for any $\alpha\le0$, $\varepsilon>0$, and $q_1,q_2\in[1,\infty]$.
\end{rema}

\begin{proof}
For~\eqref{almostequiv2}, taking $r\in[1,\infty]$ such that $1/q_1=1/r+1/q_2$, we have
\begin{multline*}
\left\|t^{-\alpha_1/\ell}\right\|Q_t f\|_{L^p(w)}\big\|_{L^{q_1}\left(0,1;t^{-1}dt\right)}\\
\le\big\|t^{(\alpha_2-\alpha_1)/\ell}\big\|_{L^r\left(0,1;\,t^{-1}dt\right)}\left\|t^{-\alpha_2/\ell}\right \|Q_t f\|_{L^p(w)}\big\|_{L^{q_2}\left(0,1;\,t^{-1}dt\right)}
\end{multline*}
by H\"older's inequality. Since $\|t^{(\alpha_2-\alpha_1)/\ell}\|_{L^r(0,1;\,t^{-1}dt)}<\infty$, we have the result. Next we prove~\eqref{almostequiv3}. By using Definition~\ref{asmp1}$\MK$\eqref{asmp2:derivative} and Lemma~\ref{lem:weightedyoung}, we have
\begin{align*}
\|Q_tf-Q_1f\|_{L^p(w)} &\le\int_t^1\|\partial_sQ_sf\|_{L^p(w)}ds =\int_t^1\|(\partial_sQ)_{s/2}Q_{s/2}f\|_{L^p(w)}ds\\
&\lesssim \int_t^1\|Q_{s/2}f\|_{L^p(w)}\frac{ds}s =\int_{t/2}^{1/2}\|Q_sf\|_{L^p(w)}\frac{ds}s.
\end{align*}
Therefore,
\begin{align*}
t^{-\alpha/\ell}\|Q_tf\|_{L^p(w)} &\lesssim t^{-\alpha/\ell}\bigg(\|Q_1f\|_{L^p(w)}+\int_{t/2}^1\|Q_sf\|_{L^p(w)}\frac{ds}s\bigg)\\
&\lesssim\|Q_1f\|_{L^p(w)}+\int_{t/2}^1s^{-\alpha/\ell}\|Q_sf\|_{L^p(w)}\frac{ds}s
\le\|f\|_{B_{p,1}^{\alpha,Q}(w)}.
\end{align*}
By taking the supremum over $t\in(0,1]$, we have the result.
\end{proof}

The following result is an analogue of the classical Besov embedding.

\begin{prop}
Let $\alpha\le0$, $p,q,r\in[1,\infty]$, and $r\ge p$. For any $f\in L^p(w)$, we have the inequality
\[
\|f\|_{B_{r,q}^{\alpha-|\mfs|\left(\frac1p-\frac1r\right),Q}(w)}\lesssim\|f\|_{B_{p,q}^{\alpha,Q}(w)}.
\]
\end{prop}

\begin{proof}
Since $\|Q_tf\|_{L^r(w)}=\|Q_{t/2}(Q_{t/2}f)\|_{L^r(w)}\lesssim t^{-\frac{|\mfs|}\ell(\frac1p-\frac1r)}\|Q_{t/2}f\|_{L^p(w)}$ by\linebreak Lemma~\ref{lem:weightedyoung}, the result follows from the definition of norms.
\end{proof}

We have the hierarchy between Besov spaces with different parameters $\alpha$.

\begin{prop}\label{prop:abinjection}
Let $\alpha_1<\alpha_2\le0$ and $p\in[1,\infty]$. The identity $\iota_{\alpha_1}:L_c^p(w)\hookrightarrow B_p^{\alpha_1,Q}(w)$ is uniquely extended to the continuous injection $\iota_{\alpha_1}^{\alpha_2}:B_p^{\alpha_2,Q}(w)\hookrightarrow B_p^{\alpha_1,Q}(w)$.
\end{prop}

\begin{proof}
We prove only the injectivity. Note that, for any $\alpha\le0$, the operator $Q_t:L_c^p(w)\to L_c^p(w)$ is continuously extended to the operator $Q_t^\alpha:B_p^{\alpha,Q}(w)\to L_c^p(w)$ and it holds that
\begin{align}\label{2:eq:besovnormex}
\|f\|_{B_p^{\alpha,Q}(w)}=\|Q_1^\alpha f\|_{L^p(w)}+\sup_{0\,<\,t\,\le\,1}t^{-\alpha/\ell}\|Q_t^\alpha f\|_{L^p(w)}
\end{align}
for any $f\in B_p^{\alpha,Q}(w)$. Let $f\in B_p^{\alpha_2,Q}(w)$ be such that $\iota_{\alpha_1}^{\alpha_2} f=0$ in $B_p^{\alpha_1,Q}(w)$. Taking a sequence $\{f_n\}\subset L_c^p(w)$ such that $f_n\to f$ in $B_p^{\alpha_2,Q}(w)$, we have $f_n=\iota_{\alpha_1}^{\alpha_2}f_n\to\iota_{\alpha_1}^{\alpha_2}f$ in $B_p^{\alpha_1,Q}(w)$ by the continuity. By the continuity of $Q_t^{\alpha_i}$ ($i\in\{1,2\}$), we have
\[
Q_t^{\alpha_2}f=\lim_{n\to\infty}Q_tf_n=Q_t^{\alpha_1}\left(\iota_{\alpha_1}^{\alpha_2}f\right)=0
\]
in $L^p(w)$ for any $t\in(0,1]$. By the identity~\eqref{2:eq:besovnormex}, we have $\|f\|_{B_p^{\alpha_2,Q}(w)}=0$.
\end{proof}

The extensions $\{Q_t^\alpha\}_{0\,<\,t\,\le\,1}$ obtained in the above proof are compatible in the sense that $Q_t^{\alpha_1}\circ\iota_{\alpha_1}^{\alpha_2}=Q_t^{\alpha_2}$. Because of this, we can omit the letter $\alpha$ and use the notation $Q_t$ to mean its extension $Q_t^\alpha$ regardless of its domain. We close this section with the continuity of $Q_t$ with respect to $t$ in Besov norms.

\begin{lemm}\label{lem:timecontinuity}
Let $\alpha\le0$ and $p\in[1,\infty]$. There exists a constant $C>0$ such that, for any $f\in B_p^{\alpha,Q}(w)$, $t\in(0,1]$, and $\varepsilon\in[0,\ell]$, we have
\[
\|(Q_t-\id)f\|_{B_p^{\alpha-\varepsilon,Q}(w)}
\le C\, t^{\varepsilon/\ell}\|f\|_{B_p^{\alpha,Q}(w)}.
\]
\end{lemm}

\begin{proof}
Similarly to the proof of Proposition~\ref{prop:almostequiv}, we have for any $s,t\in(0,1]$,
\begin{align*}
\|Q_s(Q_t-\id)f\|_{L^p(w)} &=\|(Q_{t+s}-Q_s)f\|_{L^p(w)}\le\int_s^{t+s}\|\partial_rQ_rf\|_{L^p(w)}dr\\
&\lesssim\int_s^{t+s}\|Q_{r/2}f\|_{L^p(w)}\frac{dr}r
\le\|f\|_{B_p^{\alpha,Q}(w)}\int_s^{t+s}r^{\alpha/\ell-1}dr.
\end{align*}
Since $\int_s^{t+s}r^{\alpha/\ell-1}dr\lesssim (ts^{\alpha/\ell-1})\wedge s^{\alpha/\ell}$, we have the result by an interpolation.
\end{proof}
\section{Basic notions of regularity-integrability structures}\label{section:RI}

In this section, we extend the original definitions of regularity structures, models, and modelled distributions in~\cite{Hai14} by taking integrability exponents into account.

\subsection{Regularity-integrability structures}

While the label set of the regularity structure is a set of real numbers, our label set is a subset of $\bbR\times[1,\infty]$. We denote generic elements of $\bbR\times[1,\infty]$ by bold letters $\bfa,\bfb,\bfc$, and so on. For each element $\bfa=(\alpha,p)\in\bbR\times[1,\infty]$, we write $\alpha=r(\bfa)$ and $p=i(\bfa)$, where the letters ``$r$'' and ``$i$'' mean ``regularity'' and ``integrability'', respectively. We define a partial order $\preceq$ and a strict partial order $\prec$ of the set $\bbR\times[1,\infty]$ by
\begin{align*}
\bfb\preceq\bfa&\;\overset{\deff}
{\Longleftrightarrow}\; r(\bfb)\le r(\bfa),\ i(\bfb)\ge i(\bfa),\\
\bfb\prec\bfa&\;\overset{\deff}{\Longleftrightarrow}\; r(\bfb)<r(\bfa),\ i(\bfb)\ge i(\bfa).
\end{align*}
Note that $i(\bfb)$ may be equal to $i(\bfa)$ even for the latter case. For any $\bfb\preceq\bfa$, we define the element $\bfaob\in\bbR\times[1,\infty]$ by
\[
\bfaob:=\Bigg(r(\bfa)-r(\bfb),\frac1{\frac1{i(\bfa)}-\frac1{i(\bfb)}}\Bigg),
\]
where $1/\infty:=0$ and $1/0:=\infty$. In what follows, the relations $r(\bfa)=r(\bfa\ominus\bfb)+r(\bfb)$ and $1/{i(\bfa)}=1/{i(\bfaob)}+1/{i(\bfb)}$ are important.

\begin{defi}
A \emph{regularity-integrability structure} $\scT=(\bfA,\bfT,\bfG)$ consists of the following objects.
\begin{enumerate}
\item (Index set) $\bfA$ is a subset of $\bbR\times[1,\infty]$ such that, for every $\bfa\in\bbR\times[1,\infty]$, the subset $\{\bfb\in \bfA\, ;\, \bfb\prec\bfa\}$ is finite.
\item (Model space) $\bfT=\bigoplus_{{\bfa}\,\in\,\bfA}\bfT_{\bfa}$ is an algebraic sum of Banach spaces $(\bfT_{\bfa},\|\cdot\|_{\bfa})$.
\item (Structure group) $\bfG$ is a group of continuous linear operators on $\bfT$ such that, for any $\Gamma\in \bfG$ and $\bfa\in\bfA$,
\[
(\Gamma-\id)\bfT_{\bfa}\subset\bigoplus_{{\bfb}\,\in \,\bfA,\, {\bfb}\,\prec \,{\bfa}}\bfT_{\bfb}.
\]
\end{enumerate}
A \emph{regularity} of $\scT$ is $\alpha_0\in\bbR$ such that $(\alpha_0,\infty)\preceq\bfa$ for any $\bfa\in \bfA$. For any $\bfa\in\bfA$, we denote by $P_{\bfa}:\bfT\to\bfT_{\bfa}$ a canonical projection and write
\[
\|\tau\|_{\bfa}:=\|P_{\bfa}\tau\|_{\bfa},\qquad\tau\in \bfT
\]
by abuse of notation.
\end{defi}

Obviously, the regularity structure is a particular case such that $\bfA\subset\bbR\times\{\infty\}$.

\subsection{Models}

We define the space of Besov type models on the fixed regularity-integrability structure $\scT=(\bfA,\bfT,\bfG)$. For any measurable functions $f$ on $\bbR^d$ taking values in a Banach space $(X,\|\cdot\|_X)$, we use the notation
\[
\|f\|_{L^p(w;\,X)}:=\big\|\|f(x)\|_X\big\|_{L_x^p(w)}
\]
for simplicity. For two Banach spaces $X$ and $Y$, we denote by $\mcL(X,Y)$ the Banach space of all continuous linear operators $X\to Y$.

\begin{defi}
Let $w$ be a $G$-controlled weight. A \emph{smooth model} $M=(\Pi,\Gamma)$ is a pair of two families of continuous linear operators $\Pi=\{\Pi_x:\bfT\to C(\bbR^d)\}_{x\,\in\,\bbR^d}$ and $\Gamma=\{\Gamma_{xy}\}_{x,y\,\in\,\bbR^d}\subset \bfG$ with the following properties.
\begin{enumerate}
\item (Algebraic conditions) $\Pi_x\Gamma_{xy}=\Pi_y$, $\Gamma_{xx}=\id$, and $\Gamma_{xy}\Gamma_{yz}=\Gamma_{xz}$ for any $x,y,z\in\bbR^d$.
\item (Analytic conditions) For any $\bfc\in\bbR\times[1,\infty]$,
\begin{align*}
\|\Pi\|_{\bfc,w}&:=\max_{\bfa\,\in\,\bfA,\: \bfa\,\prec\,\bfc}\;
\sup_{0\,<\,t\,\le\,1}
\Big(t^{-r(\bfa)/\ell}\big\|Q_t\big(x,\Pi_x(\cdot)\big)\big\|_{L_x^{i(\bfa)}\left(w;\,\bfT_{\bfa}^*\right)}\Big)
\\
&=\max_{\bfa\,\in\,\bfA,\, \bfa\,\prec\,\bfc}\,
\sup_{0\,<\,t\,\le\,1}
\left(t^{-r(\bfa)/\ell}
\left\|
\sup_{\tau\,\in\,\bfT_{\bfa}\setminus\{0\}}
\frac{|Q_t(x,\Pi_x\tau)|}{\|\tau\|_{\bfa}}
\right\|_{L_x^{i(\bfa)}(w)}\right)<\infty
\end{align*}
and
\begin{align*}
&\|\Gamma\|_{\bfc,w}:=\max_{\substack{\bfa,\bfb\,\in\,\bfA\\ \bfb\,\prec\,\bfa\,\prec\,\bfc}}\,
\sup_{h\,\in\,\bbR^d\setminus\{0\}}
\frac{\|\Gamma_{(x+h)x}\|_{L_x^{i(\bfaob)}\left(w;\,\mcL(\bfT_{\bfa},\bfT_{\bfb})\right)}}{\wco(h)\|h\|_\mfs^{r(\bfaob)}}\\
&=\max_{\substack{\bfa,\bfb\,\in\,\bfA\\ \bfb\,\prec\,\bfa\,\prec\,\bfc}}\,
\sup_{h\,\in\,\bbR^d\setminus\{0\}}\left(
\frac{1}{\wco(h)\|h\|_\mfs^{r(\bfaob)}}
\left\|
\sup_{\tau\,\in\,\bfT_{\bfa}\setminus\{0\}}
\frac{\|\Gamma_{(x+h)x}\tau\|_{\bfb}}{\|\tau\|_{\bfa}}
\right\|_{L_x^{i(\bfaob)}(w)}\right)<\infty.
\end{align*}
\end{enumerate}
We write $\tri M\tri_{\bfc,w}:=\|\Pi\|_{\bfc,w}+\|\Gamma\|_{\bfc,w}$. In addition, for any two smooth models $M^{(i)}=(\Pi^{(i)},\Gamma^{(i)})$ with $i\in\{1,2\}$, we define the pseudo-metrics
\[
\bigtri M^{(1)};M^{(2)}\bigtri_{\bfc,w}:=\left\|\Pi^{(1)}-\Pi^{(2)}\right\|_{\bfc,w}+\left\|\Gamma^{(1)}-\Gamma^{(2)}\right\|_{\bfc,w}
\]
by replacing $\Pi$ and $\Gamma$ above with $\Pi^{(1)}-\Pi^{(2)}$ and $\Gamma^{(1)}-\Gamma^{(2)}$ respectively. Finally, we define the space $\scM_w(\scT)$ as the completion of the set of all smooth models, under the pseudo-metrics $\tri\cdot;\cdot\tri_{\bfc,w}$ for all $\bfc\in\bbR\times[1,\infty]$. We call each element of $\scM_w(\scT)$ a \emph{model} for $\scT$. We still use the notation $M=(\Pi,\Gamma)$ to denote a generic model.
\end{defi}

Recall that $i(\bfaob)=\infty$ if $\bfa,\bfb\in\bbR\times\{\infty\}$. Therefore, in the regularity structure case $\bfA\subset\bbR\times\{\infty\}$, the above definition coincides with the original definition of models~\cite[Definition~2.17]{Hai14} if we ignore the difference between local and global bounds.

%\medskip

It is a subtle question in which space the operator ``$\Pi_x$'' takes values for general $M=(\Pi,\Gamma)\in\scM_w(\scT)$. Under some additional assumptions on weights, we can regard $\Pi_x$ as a continuous linear operator from $\bfT$ to a Besov space.

\begin{prop}\label{prop:whatisPixtau}
Let $\alpha_0\le0$ be a regularity of $\scT$. Assume that there exist two $G$-controlled weights $w_1$ and $w_2$ such that
\[
\sup_{x\,\in\,\bbR^d}\big\{\|x\|_\mfs^n\,\wco(x)w_1(x)\big\} +\sup_{x\,\in\,\bbR^d}\big\{\|x\|_\mfs^n\,\wco_1(x)w_2(x)\big\} <\infty
\]
for any $n\ge0$, and that $ww_1$ and $ww_2$ are also $G$-controlled. Then for almost every $x\in\bbR^d$, the map $\Pi_x$ is well-defined as a continuous linear operator from $\bfT_{\bfa}$ to $B_{i(\bfa),1}^{\alpha,Q}(ww_1)$ for any $\bfa\in\bfA$ and any $\alpha<\alpha_0$. More precisely, for any $\bfc\in\bbR\times[1,\infty]$ such that $\bfa\prec\bfc$ we have
\[
\|\Pi_x\|_{L_x^{i(\bfa)}\big(ww_2;\,\mcL\big(\bfT_\bfa,B_{i(\bfa),1}^{\alpha,Q}(ww_1)\big)\big)}
\lesssim\|\Pi\|_{\bfc,w}(1+\|\Gamma\|_{\bfc,w}).
\]
\end{prop}

\begin{proof}
By the density argument, it is sufficient to show the inequality for smooth models. By the algebraic relations and Lemmas~\ref{lem:weightedholder} and~\ref{lem:weightedtranslation},
\begingroup
\allowdisplaybreaks
\begin{multline*}
\big\|Q_t\big(y,\Pi_x(\cdot)\big)\big\|_{\mcL\left(\bfT_\bfa,L_y^{i(\bfa)}(ww_1)\right)}\\
\begin{aligned}
&=\sup_{\tau\,\in\,\bfT_\bfa\setminus\{0\}}\frac{\|Q_t(y,\Pi_y\Gamma_{yx}\tau)\|_{L_y^{i(\bfa)}(ww_1)}}{\|\tau\|_\bfa}\\
&\le\sum_{\bfb\,\preceq\,\bfa}\left\|\big\|Q_t\big(y,\Pi_y(\cdot)\big)\big\|_{\bfT_{\bfb}^*}
\sup_{\tau\,\in\,\bfT_\bfa\setminus\{0\}}\frac{\|\Gamma_{yx}\tau\|_{\bfb}}{\|\tau\|_\bfa}\right\|_{L_y^{i(\bfa)}(ww_1)}\\
&\le \sum_{\bfb\,\preceq\,\bfa}\big\|Q_t\big(y,\Pi_y(\cdot)\big)\big\|_{L_y^{i(\bfb)}\left(w;\,\bfT_\bfb^*\right)}
\big\|\|\Gamma_{yx}\|_{\mcL\left(\bfT_\bfa,\bfT_\bfb\right)}\big\|_{L_y^{i(\bfa\ominus\bfb)}(w_1)}\\
&\le \|\Pi\|_{\bfc,w}\sum_{\bfb\,\preceq\,\bfa}t^{r(\bfb)/\ell}\wco_1(x)
\big\|\|\Gamma_{(x+y)x}\|_{\mcL(\bfT_\bfa,\bfT_\bfb)}\big\|_{L_y^{i(\bfa\ominus\bfb)}(w_1)}.
\end{aligned}
\end{multline*}
\endgroup
By H\"older's inequality and Fubini's theorem, we have
\begingroup
\allowdisplaybreaks
\begin{multline*}
\Big\|\big\|Q_t\left(y,\Pi_x(\cdot)\right)\big\|_{\mcL(\bfT_\bfa,L_y^{i(\bfa)}(ww_1))}\Big\|_{L_x^{i(\bfa)}(ww_2)}\\
\begin{aligned}
&\le \|\Pi\|_{\bfc,w}\sum_{\bfb\,\preceq\,\bfa}t^{r(\bfb)/\ell}\|\wco_1\|_{L^{i(\bfb)}(w_2)}
\Big\|\big\|\|\Gamma_{(x+y)x}\|_{\mcL(\bfT_\bfa,\bfT_\bfb)}\big\|_{L_y^{i(\bfaob)}(w_1)}\Big\|_{L_x^{i(\bfaob)}(w)}\\
&=\|\Pi\|_{\bfc,w}\sum_{\bfb\,\preceq\,\bfa}t^{r(\bfb)/\ell}\|\wco_1\|_{L^{i(\bfb)}(w_2)}
\Big\|\big\|\|\Gamma_{(x+y)x}\|_{\mcL(\bfT_\bfa,\bfT_\bfb)}\big\|_{L_x^{i(\bfaob)}(w)}\Big\|_{L_y^{i(\bfaob)}(w_1)}\\
&\lesssim\|\Pi\|_{\bfc,w}(1+\|\Gamma\|_{\bfc,w})\sum_{\bfb\,\preceq\,\bfa}t^{r(\bfb)/\ell}\big\|1+\wco(y)\|y\|_\mfs^{r(\bfaob)}\big\|_{L_y^{i(\bfaob)}(w_1)}\\
&\lesssim\|\Pi\|_{\bfc,w}(1+\|\Gamma\|_{\bfc,w})\, t^{\alpha_0/\ell}
\end{aligned}
\end{multline*}
\endgroup
for $t\in(0,1]$. Note that $w^*\ge1$ and $\|1\|_{L^p(w_1)}<\infty$ for any $p\in[1,\infty]$ by an assumption on $w_1$. Finally, by the definition of Besov norms,
\begin{align*}
&\Big\|\|\Pi_x\|_{\mcL\big(\bfT_\bfa,B_{i(\bfa),1}^{\alpha,Q}(ww_1)\big)}\Big\|_{L_x^{i(\bfa)}(ww_2)}\\
&\le\bigg\|
\big\|Q_1\big(y,\Pi_x(\cdot)\big)\big\|_{\mcL\left(\bfT_\bfa,L_y^{i(\bfa)}(ww_1)\right)}\\
&\mkern150mu+\int_0^1t^{-\alpha/\ell}\big\|Q_t\big(y,\Pi_x(\cdot)\big)\big\|_{\mcL\left(\bfT_\bfa,L_y^{i(\bfa)}(ww_1)\right)}\frac{dt}t
\bigg\|_{L_x^{i(\bfa)}(ww_2)}\\
&\lesssim\|\Pi\|_{\bfc,w}(1+\|\Gamma\|_{\bfc,w})\bigg(1+\int_0^1t^{(\alpha_0-\alpha)/\ell}\frac{dt}t\bigg)
\lesssim\|\Pi\|_{\bfc,w}(1+\|\Gamma\|_{\bfc,w}).\qedhere
\end{align*}
\end{proof}

\begin{rema}\label{rem:whatisknownforZ}
Without additional weights as above, we only know that $Q_t(x-h,\Pi_x(\cdot))$ is defined for any $h\in\bbR^d$ as an element of $L_x^{i(\bfa)}(w^2;\bfT_\bfa^*)$. Indeed, for any smooth model and for any $\bfa\prec\bfc$, by Lemmas~\ref{lem:weightedholder} and~\ref{lem:weightedtranslation} we have
\begin{align}\label{eq:Q(x-h,Pix)}
\begin{aligned}
\big\|Q_t\big(x-h,&\Pi_x(\cdot)\big)\big\|_{L_x^{i(\bfa)}\left(w^2;\,\bfT_\bfa^*\right)} =\big\|Q_t\big(x-h,\Pi_{x-h}\Gamma_{(x-h)x}(\cdot)\big)\big\|_{L_x^{i(\bfa)}\left(w^2;\,\bfT_\bfa^*\right)}\\
&\le\sum_{\bfb\,\preceq\,\bfa}\Big\|\big\|Q_t\big(x-h,\Pi_{x-h}(\cdot)\big)\big\|_{\bfT_\bfb^*}\|\Gamma_{(x-h)x}\|_{\mcL(\bfT_\bfa,\bfT_\bfb)}\Big\|_{L_x^{i(\bfa)}(w^2)}\\
&\le \wco(h)\sum_{\bfb\,\preceq\,\bfa}\big\|Q_t\big(x,\Pi_x(\cdot)\big)\big\|_{L_x^{i(\bfb)}\left(w;\,\bfT_\bfb^*\right)}\|\Gamma_{(x-h)x}\|_{L_x^{i(\bfaob)}(w;\,\mcL(\bfT_\bfa,\bfT_\bfb))}\\
&\le \big(\wco(h)\big)^2\|\Pi\|_{\bfc,w}(1+\|\Gamma\|_{\bfc,w})\sum_{\bfb\,\preceq\,\bfa}t^{r(\bfb)/\ell}\|h\|_\mfs^{r(\bfaob)}<\infty.
\end{aligned}
\end{align}
Moreover, by the density argument we also have the semigroup property
\[
Q_t(x,\Pi_x\tau)=\int_{\bbR^d}Q_{t-s}(x,x-h)Q_s(x-h,\Pi_x\tau)dh,\qquad 0<s<t
\]
for any models. These properties are used to prove the reconstruction theorem.
\end{rema}

\subsection{Modelled distributions}

We close this section with the definition of Besov type modelled distributions and their reconstructions. We fix two $G$-controlled weights $w$ and $v$ such that $wv$ is also $G$-controlled.


\begin{defi}\label{def:besovMD}
Let $M=(\Pi,\Gamma)\in\scM_{w}(\scT)$. For any $\bfc\in\bbR\times[1,\infty]$, we define $\mcD_{v}^\bfc(\Gamma)$ as the space of all functions $f:\bbR^d\to \bfT_{\prec \bfc}:=\bigoplus_{\bfa\,\in\,\bfA,\, \bfa\,\prec \,\bfc}\bfT_\bfa$ such that
\begin{align*}
\lp f\rp_{\bfc,v}&:=\max_{\bfa\,\prec \,\bfc}\|f\|_{L^{i(\bfcoa)}(v;\,\bfT_\bfa)}<\infty,\\
\| f\|_{\bfc,v}^\Gamma&:=\max_{\bfa\prec\bfc}
\sup_{h\,\in\,\bbR^d\setminus\{0\}}
\frac{\left\|\Delta_{x;\,h}^\Gamma f\right\|_{L_x^{i(\bfcoa)}(v;\,\bfT_\bfa)}}{\vco(h)\|h\|_\mfs^{r(\bfcoa)}}<\infty,
\end{align*}
where $\Delta_{x;h}^\Gamma f:=f(x-h)-\Gamma_{(x-h)x}f(x)$. We write $\tri f\tri_{\bfc,v}^\Gamma:=\lp f\rp_{\bfc,v}+\|f\|_{\bfc,v}^\Gamma$. We call each element of $\mcD_v^{\bfc}(\Gamma)$ a \emph{modelled distribution}.

In addition, for any two models $M^{(i)}=(\Pi^{(i)},\Gamma^{(i)})\in\scM_{w}(\scT)$ and modelled distributions $f^{(i)}\in\mcD_{v}^\bfc(\Gamma^{(i)})$ with $i\in\{1,2\}$, we define
\begin{align*}
\bigtri f^{(1)};f^{(2)}\bigtri_{\bfc,v}^{\Gamma^{(1)};\,\Gamma^{(2)}} &:=\lbp f^{(1)}-f^{(2)}\rbp_{\bfc,v}+\left\|f^{(1)};f^{(2)}\right\|_{\bfc,v}^{\Gamma^{(1)};\,\Gamma^{(2)}}
\\
\intertext{by}
\lbp f^{(1)}-f^{(2)}\rbp_{\bfc,v}&:=\max_{\bfa\,\prec \,\bfc}\left\|f^{(1)}-f^{(2)}\right\|_{L^{i(\bfcoa)}\left(v;\,\bfT_\bfa\right)},\\
\left\| f^{(1)};f^{(2)}\right\|_{\bfc,v}^{\Gamma^{(1)};\,\Gamma^{(2)}}&:=\max_{\bfa\,\prec\, \bfc}
\sup_{h\,\in\,\bbR^d\setminus\{0\}}
\frac{\left\|\Delta_{x;\,h}^{\Gamma^{(1)}} f^{(1)}-\Delta_{x;\,h}^{\Gamma^{(2)}}f^{(2)}\right\|_{L_x^{i(\bfcoa)}\left(v;\bfT_\bfa\right)}}{\vco(h)\|h\|_\mfs^{r(\bfcoa)}}.
\end{align*}
We omit the symbol ``$\,\Gamma^{(1)};\Gamma^{(2)}$'' below for simplicity.
\end{defi}

\begin{defi}\label{def:besovreconst}
Let $\alpha_0\le0$ be a regularity of $\scT$ and let $\bfc\in\bbR\times[1,\infty]$. For any $M=(\Pi,\Gamma)\in\scM_{w}(\scT)$ and $f\in\mcD_{v}^\bfc(\Gamma)$, we call any $\Lambda\in B_{i(\bfc)}^{\alpha_0,Q}(wv)$ satisfying
\begin{align*}
\lb\Lambda\rb_{\bfc,wv}^{\Pi,f}:=\sup_{0\,<\,t\,\le\,1}t^{-r(\bfc)/\ell}\big\|Q_t(x,\Lambda)-Q_t\big(x,\Pi_xf(x)\big)\big\|_{L_x^{i(\bfc)}(wv)}<\infty
\end{align*}
a \emph{reconstruction} of $f$ for $M$. Furthermore, for any models $M^{(i)}=(\Pi^{(i)},\Gamma^{(i)})\in\scM_{w}(\scT)$, modelled distributions $f^{(i)}\in\mcD_{v}^e(\Gamma^{(i)})$, and any reconstructions $\Lambda^{(i)}\in B_{i(\bfc)}^{\alpha_0,Q}(wv)$ of $f^{(i)}$ for $M^{(i)}$ with $i\in\{1,2\}$, we define
\begin{align*}
\lbb\Lambda^{(1)};\Lambda^{(2)}\rbb_{\bfc,wv}^{\Pi^{(1)},\,f^{(1)};\,\Pi^{(2)},\,f^{(2)}} :=&\sup_{0\,<\,t\,\le\,1}t^{-r(\bfc)/\ell}\Big\|\big\{Q_t\big(x,\Lambda^{(1)}\big)-Q_t\big(x,\Pi_x^{(1)}f^{(1)}(x)\big)\big\}\\
&\qquad-\big\{Q_t\big(x,\Lambda^{(2)}\big)-Q_t\big(x,\Pi_x^{(2)}f^{(2)}(x)\big)\big\}\Big\|_{L_x^{i(\bfc)}(wv)}.
\end{align*}
We also omit the symbol ``$\,\Pi^{(1)},f^{(1)};\Pi^{(2)},f^{(2)}$'' below for simplicity.
\end{defi}

\begin{rema}\label{rem:L-Pixfx}
It seems more natural to write ``$\, Q_t(x,\Lambda-\Pi_xf(x))$'', but we split it into two terms here to avoid the subtle question of what ``$\, \Pi_xf(x)$'' is (see Proposition~\ref{prop:whatisPixtau}). Since $Q_t(x,\Pi_x(\cdot))$ is well-defined as an element of $L_x^{i(\bfa)}(w;\bfT_\bfa^*)$, we can define the quantity $Q_t(x,\Pi_xf(x))$ by inserting $f(x)$ into the operator $Q_t(x,\Pi_x(\cdot))$. See also the calculations at the beginning of the proof of Theorem~\ref{thm:besovreconstruction}. We can also define $Q_t(x-h,\Pi_xf(x))$ for any $h\in\bbR^d$ by Remark~\ref{rem:whatisknownforZ}.
\end{rema}

\section{Reconstruction theorem}\label{section:reconstruction}


In this section, we fix a regularity-integrability structure $\scT=(\bfA,\bfT,\bfG)$ of regularity $\alpha_0\le0$, and also fix $G$-controlled weights $w$ and $v$ such that $wv$ is also $G$-controlled.
\begin{theo}\label{thm:besovreconstruction}
Let $\bfc\in(0,\infty)\times[1,\infty]$. Then for any $M=(\Pi,\Gamma)\in\scM_{w}(\scT)$ and $f\in\mcD_{v}^\bfc(\Gamma)$, there exists a unique reconstruction $\mcR f$ of $f$ for $M$ and it holds that
\begin{align}\label{besovreconstruction}
\|\mcR f\|_{B_{i(\bfc)}^{\alpha_0,Q}(wv)}&\lesssim\|\Pi\|_{\bfc,w}\tri f\tri_{\bfc,v}^\Gamma,\\\label{besovreconstructioncharacterize}
\lb\mcR f\rb_{\bfc,wv}^{\Pi,f}&\lesssim \|\Pi\|_{\bfc,w}\|f\|_{\bfc,v}^\Gamma.
\end{align}
\end{theo}
\begin{rema}
The original reconstruction theorem~\cite[Theorem~3.10]{Hai14} was extended to different types of norms~\cite{BL22,HL17, HR20}. Hairer and Labbé~\cite{HL17} proved a reconstruction theorem for $B_{\infty,\infty}$ type models and $B_{p,q}$ type modelled distributions. Their result is a special case of Theorem~\ref{thm:besovreconstruction} if we ignore $q$-exponents. Broux and Lee~\cite{BL22} proved Besov reconstruction theorem at the level of ``coherent germ'', which was the notion introduced by Caravenna and Zambotti~\cite{CZ20} to reformulate the reconstruction theorem without using regularity structures. As seen in the following proof, our situation is contained in~\cite{BL22} as a special case because $\{F_x:=\Pi_xf(x)\}_{x\,\in\,\bbR^d}$ is a coherent germ. However, the detailed regularity-integrability structure is effectively used in the paper~\cite{BPHZ}. As for the different norm, Hensel and Rosati~\cite{HR20} proved Triebel--Lizorkin type reconstruction theorem for $F_{\infty,\infty}$ type models and $F_{p,q}$ type modelled distributions.
\end{rema}

\begin{proof}
It is sufficient to show the bounds~\eqref{besovreconstruction} and~\eqref{besovreconstructioncharacterize} for a single model and modelled distribution. The proofs of the local Lipschitz estimates are simple modifications. For $t>0$ and $0<s\le t\wedge1$, we define the functions
\begin{align*}
\mcR_s^t f(x):=
%\left\{
\begin{alcases}
&\int_{\bbR^d}Q_{t-s}(x,y)Q_s\big(y,\Pi_yf(y)\big)dy,\quad&s<t,\\
&Q_t\big(x,\Pi_xf(x)\big),&s=t.
\end{alcases}
\end{align*}
By Lemma~\ref{lem:weightedholder}, we have
\begin{align*}
\big\|Q_s\big(y,\Pi_yf(y)\big)\big\|_{L_y^{i(\bfc)}(wv)} &\le\sum_{\bfa\,\prec \,\bfc}\Big\|\big\|Q_s\big(y,\Pi_y(\cdot)\big)\big\|_{\bfT_\bfa^*}\|P_\bfa f(y)\|_{\bfT_\bfa}\Big\|_{L_y^{i(\bfc)}(wv)}\\
&\le\sum_{\bfa\,\prec \,\bfc}\big\|Q_s\big(y,\Pi_y(\cdot)\big)\big\|_{L_y^{i(\bfa)}(w;\,\bfT_\bfa^*)}\|f\|_{L^{i(\bfcoa)}(v;\,\bfT_\bfa)}\\
&\le \|\Pi\|_{\bfc,w}\lp f\rp_{\bfc,v}\sum_{\bfa\,\prec \,\bfc}s^{r(\bfa)/\ell}
\lesssim\|\Pi\|_{\bfc,w}\lp f\rp_{\bfc,v}\, s^{\alpha_0/\ell}.
\end{align*}
Hence by Proposition~\ref{prop:operatorQt}, we have $\mcR_s^tf\in L_c^{i(\bfc)}(wv)$ and
\[
\left\|\mcR_s^tf\right\|_{L^{i(\bfc)}(wv)}\lesssim\|\Pi\|_{\bfc,w}\lp f\rp_{\bfc,v}\, s^{\alpha_0/\ell}.
\]
We separate the proof into five steps.

%\medskip

\paragraph*{(1) Coherence property.}
Set $F_x:=\Pi_xf(x)$. (This is an abuse of notation as mentioned in Remark~\ref{rem:L-Pixfx}, but it does not cause a serious problem because it always appears in the form $Q_t(x-h,F_x)$.) By Lemmas~\ref{lem:weightedholder} and~\ref{lem:weightedtranslation}, we have
\begin{multline*}
\big\|Q_t(x-h,F_x-F_{x-h})\big\|_{L_x^{i(\bfc)}(wv)}
\\
\begin{aligned}
&=\big\|Q_t\big(x-h,\Pi_{x-h}\big\{\Gamma_{(x-h)x}f(x)-f(x-h)\big\}\big)\big\|_{L_x^{i(\bfc)}(wv)}\\
&\le\sum_{\bfa\,\prec \,\bfc}\big\|Q_t\big(x-h,\Pi_{x-h}(\cdot)\big)\big\|_{L_x^{i(\bfa)}\left(w;\,\bfT_\bfa^*\right)}\left\|\Delta_{x;h}^\Gamma f\right\|_{L_x^{i(\bfcoa)}(v;\,\bfT_\bfa)}
\\
&\le \|\Pi\|_{\bfc,w}\|f\|_{\bfc,v}^\Gamma\,\wco(h)\sum_{\bfa\,\prec \,\bfc}t^{r(\bfa)/\ell}\vco(h)\|h\|_{\mfs}^{r(\bfcoa)}.
\end{aligned}
\end{multline*}

\paragraph*{(2) Convergence as $s\downarrow0$.}
By the semigroup property of $\{Q_t\}_{t\,>\,0}$, for $0<u<s<t\wedge1$,
\begin{align*}
\left|\mcR_s^tf(x)-\mcR_u^tf(x)\right| &=\bigg|\int_{(\bbR^d)^2}Q_{t-s}(x,y)Q_{s-u}(y,y-h)Q_u(y-h,F_y-F_{y-h})dydh\bigg|\\
&\lesssim\int_{(\bbR^d)^2}G_{t-s}(x-y)G_{s-u}(h)\big|Q_u(y-h,F_y-F_{y-h})\big|dydh.
\end{align*}
By applying Lemma~\ref{lem:weightedyoung} to the convolution with respect to $y$,
\begin{align*}
\left\|\mcR_u^tf-\mcR_s^tf\right\|_{L^{i(\bfc)}(wv)} &\lesssim\int_{\bbR^d}G_{s-u}(h)\big\|Q_u(y-h,F_{y}-F_{y-h})\big\|_{L_y^{i(\bfc)}(wv)}dh\\
&\le\|\Pi\|_{\bfc,w}\|f\|_{\bfc,v}^\Gamma\sum_{\bfa\prec\bfc}u^{r(\bfa)/\ell}\int_{\bbR^d}\|h\|_{\mfs}^{r(\bfcoa)}(\wco\vco)(h)G_{s-u}(h)dh\\
&\lesssim
\|\Pi\|_{\bfc,w}\|f\|_{\bfc,v}^\Gamma\sum_{\bfa\prec\bfc}u^{r(\bfa)/\ell}(s-u)^{r(\bfcoa)/\ell}.
\end{align*}
In the last inequality, we use the property~\eqref{weightintegrable} for $\wco\vco$. Since $r(\bfa)+r(\bfcoa)=r(\bfc)$, we have for any $u\in[s/2,s)$,
\begin{align}\label{eq:besovreconstructioncauchy}
\left\|\mcR_u^tf-\mcR_s^tf\right\|_{L^{i(\bfc)}(wv)}\lesssim\|\Pi\|_{\bfc,w}\|f\|_{\bfc,v}^\Gamma\, s^{r(\bfc)/\ell}.
\end{align}
Also for $u\in(0,s/2)$, by taking $n\in\bbN$ such that $u\in[s/2^{n+1},s/2^n)$, we have
\begin{align*}
\left\|\mcR_u^tf-\mcR_s^tf\right\|_{L^{i(\bfc)}(wv)} &\le\sum_{m=0}^{n-1}\left\|\mcR_{s/2^m}^tf-\mcR_{s/2^{m+1}}^tf\right\|_{L^{i(\bfc)}(wv)}+\left\|\mcR_{s/2^n}^tf-\mcR_u^tf\right\|_{L^{i(\bfc)}(wv)}\\
&\lesssim\|\Pi\|_{\bfc,w}\|f\|_{\bfc,v}^\Gamma\left
\{\sum_{m=0}^{n-1}\left(\frac{s}{2^m}\right)^{r(\bfc)/\ell}+\left(\frac{s}{2^n}\right)^{r(\bfc)/\ell}\right\}\\
&\lesssim\|\Pi\|_{\bfc,w}\|f\|_{\bfc,v}^\Gamma\, s^{r(\bfc)/\ell}.
\end{align*}
Moreover, the same bound for the case $s=t\le1$ can be obtained by a similar argument. In the end, the bound~\eqref{eq:besovreconstructioncauchy} holds for any $0<u<s\le t\wedge1$. Since $r(\bfc)>0$, this implies that $\{\mcR_s^tf\}_{0\,<\,s\,\le\,t\,\wedge\,1}$ is Cauchy in $L_c^{i(\bfc)}(wv)$ as $s\downarrow0$. We denote its limit by
\[
\mcR_0^tf:=\lim_{s\downarrow0}\mcR_s^tf.
\]

\paragraph*{(3) Uniform bounds.} Combining the Cauchy property~\eqref{eq:besovreconstructioncauchy} with the initial bound
\[
\left\|\mcR_{t\wedge1}^tf\right\|_{L^{i(\bfc)}(wv)}\lesssim\|\Pi\|_{\bfc,w}\lp f\rp_{\bfc,v}(t\wedge1)^{\alpha_0/\ell}
\]
obtained at the beginning of the proof, we have
\begin{align}\label{eq:besovreconstructuniform}
\left\|\mcR_0^tf\right\|_{L^{i(\bfc)}(wv)}
\lesssim\|\Pi\|_{\bfc,w}\tri f \tri_{\bfc,v}^\Gamma(t\wedge1)^{\alpha_0/\ell}.
\end{align}
Incidentally, we have the identity
\[
Q_s\mcR_u^tf=\mcR_u^{t+s}f,\qquad 0<u\le t\wedge1,\ s>0
\]
from the semigroup property. Letting $u\downarrow0$, we have
\begin{align}\label{eq:besovreconstructsemigroupR0}
Q_s\mcR_0^tf=\mcR_0^{t+s}f,\qquad t,s>0.
\end{align}
By~\eqref{eq:besovreconstructsemigroupR0} and the bound~\eqref{eq:besovreconstructuniform}, we have the estimate of $\mcR_0^tf$ in the Besov norm
\begin{align}\label{eq:besovreconstructBesovbdd}
\left\|\mcR_0^tf\right\|_{B_{i(\bfc)}^{\alpha_0,Q}(wv)} &=\sup_{0\,<\,s\,\le\,1}s^{-\alpha_0/\ell}\left\|\mcR_0^{t+s}f\right\|_{L^{i(\bfc)}(wv)}
\lesssim\|\Pi\|_{\bfc,w}\tri f \tri_{\bfc,v}^\Gamma.
\end{align}

%\medskip

\paragraph*{(4) Convergence as $t\downarrow0$.}
By the semigroup property~\eqref{eq:besovreconstructsemigroupR0}, the uniform bound~\eqref{eq:besovreconstructBesovbdd}, and Lemma~\ref{lem:timecontinuity}, we have for any $\varepsilon\in(0,\ell]$ and any $0<s<t$,
\begin{align*}
\left\|\mcR_0^tf-\mcR_0^sf\right\|_{B_{i(\bfc)}^{\alpha_0-\varepsilon,Q}(wv)} &=\left\|(Q_{t-s}-\id)\mcR_0^sf\right\|_{B_{i(\bfc)}^{\alpha_0-\varepsilon,Q}(wv)}
\lesssim(s-r)^{\varepsilon/\ell}\|\Pi\|_{\bfc,w}\tri f \tri_{\bfc,v}^\Gamma.
\end{align*}
This implies that $\{\mcR_0^tf\}_{t\,\in\,(0,1]}$ is Cauchy in $B_{i(\bfc)}^{\alpha_0-\varepsilon,Q}(wv)$ as $t\downarrow0$. We denote its limit by
\[
\mcR f:=\lim_{t\downarrow0}\mcR_0^tf.
\]
Incidentally, by letting $t\downarrow0$ in~\eqref{eq:besovreconstructsemigroupR0}, we have
\[
Q_s\mcR f=\mcR_0^sf,\quad s>0.
\]
Combining this with the uniform bound~\eqref{eq:besovreconstructuniform}, we have that $\mcR f$ actually belongs to $B_{i(\bfc)}^{\alpha_0,Q}(wv)$ and the result~\eqref{besovreconstruction} follows. On the other hand, by letting $u\downarrow0$ and $s=t$ in~\eqref{eq:besovreconstructioncauchy}, we have the result~\eqref{besovreconstructioncharacterize}.

%\medskip

\paragraph*{(5) Uniqueness.}
Let $\Lambda,\Lambda'\in B_{i(\bfc)}^{\alpha_0,Q}(wv)$ be reconstructions of $f$ for $M$. From the definition of reconstruction, $g:=\Lambda-\Lambda'\in B_{i(\bfc)}^{\alpha_0,Q}(wv)$ satisfies
\begin{multline*}
\|Q_tg\|_{L^{i(\bfc)}(wv)}\\
\begin{aligned}
&\le\big\|Q_t(x,\Lambda)-Q_t\big(x,\Pi_xf(x)\big)\big\|_{L_x^{i(\bfc)}(wv)}+\big\|Q_t(x,\Lambda')-Q_t\big(x,\Pi_xf(x)\big)\big\|_{L_x^{i(\bfc)}(wv)}\\
&\lesssim t^{r(\bfc)/\ell}.
\end{aligned}
\end{multline*}
Then by using Lemma~\ref{lem:timecontinuity} again, we have that for any $\varepsilon\in(0,\ell]$
\begin{align*}
\|g\|_{B_{i(\bfc)}^{\alpha_0-\varepsilon,Q}(wv)}&\le\|(Q_t-\id)g\|_{B_{i(\bfc)}^{\alpha_0-\varepsilon,Q}(wv)}+\|Q_tg\|_{B_{i(\bfc)}^{\alpha_0-\varepsilon,Q}(wv)}\\
&\lesssim t^{\varepsilon/\ell}\|g\|_{B_{i(\bfc)}^{\alpha_0,Q}(wv)}+\|Q_tg\|_{L^{i(\bfc)}(wv)}
\lesssim t^{(\varepsilon\wedge r(\bfc))/\ell}.
\end{align*}
Since $r(\bfc)>0$, we have $g=0$ in $B_{i(\bfc)}^{\alpha_0-\varepsilon,Q}(wv)$ by taking the limit $t\downarrow0$. By Proposition~\ref{prop:abinjection}, this implies $g=0$ in $B_{i(\bfc)}^{\alpha_0,Q}(wv)$, and thus $\Lambda=\Lambda'$.
\end{proof}

\begin{thebibliography}{BDY12}
\bibitem[BB16]{BB16} Bailleul, Ismaël F.; Bernicot, Frédéric Heat semigroup and singular PDEs, J. Funct. Anal., Volume 270 (2016) no. 9, pp. 3344-3452.
\bibitem[BCD11]{BCD11} Bahouri, Hajer; Chemin, Jean-Yves; Danchin, Raphaël Fourier analysis and nonlinear partial differential equations, Grundlehren der Mathematischen Wissenschaften, 343, Springer, 2011.
\bibitem[BDY12]{BDY12} Bui, Huy-Qui; Duong, Xuan Thinh; Yan, Lixin Calderón reproducing formulas and new Besov spaces associated with operators, Adv. Math., Volume 229 (2012) no. 4, pp. 2449-2502.
\bibitem[BH24]{BPHZ} Bailleul, Ismaël F.; Hoshino, Masato Random models on regularity-integrability structures (2024) (https://arxiv.org/abs/2310.10202).
\bibitem[BL23]{BL22} Broux, Lucas; Lee, David Besov reconstruction, Potential Anal., Volume 59 (2023) no. 4, pp. 1875-1912.
\bibitem[CZ20]{CZ20} Caravenna, Francesco; Zambotti, Lorenzo Hairer’s reconstruction theorem without regularity structures, EMS Surv. Math. Sci., Volume 7 (2020) no. 2, pp. 207-251.
\bibitem[GL15]{GL15} Grigor’yan, Alexander; Liu, Liguang Heat kernel and Lipschitz–Besov spaces, Forum Math., Volume 27 (2015) no. 6, pp. 3567-3613.
\bibitem[HL17]{HL17} Hairer, Martin; Labbé, Cyril The reconstruction theorem in Besov spaces, J. Funct. Anal., Volume 273 (2017) no. 8, pp. 2578-2618.
\bibitem[HR20]{HR20} Hensel, Sebastian; Rosati, Tommaso Modelled distributions of Triebel–Lizorkin type, Stud. Math., Volume 252 (2020) no. 3, pp. 251-297.
\bibitem[Hai14]{Hai14} Hairer, Martin A theory of regularity structures, Invent. Math., Volume 198 (2014) no. 2, pp. 269-504.
\bibitem[MW17]{MW17} Jean-Christophe Mourrat and Hendrik Weber, Global well-posedness of the dynamic $\Phi^4$ model in the plane, Ann. Probab.45(2017), no.4,2398–2476.
\bibitem[Tri92]{Tri92} Triebel, Hans Theory of function spaces II, Monographs in Mathematics, 84, Birkhäuser, 1992.
\end{thebibliography}
\end{document}
